\title{Direct Preference Optimization: Your Language Model is Secretly a Reward Model}

\begin{abstract}
Although large-scale unsupervised language models (LMs) are capable of acquiring extensive world knowledge and some degree of reasoning through their training, exerting fine-grained control over their outputs is challenging due to the lack of supervision. Current approaches to making these models more steerable typically involve collecting human judgments that rank the relative quality of model outputs, then adjusting the unsupervised LM via fine-tuning to better align with these human preferences—often using reinforcement learning from human feedback (RLHF). However, RLHF presents considerable complexity and instability, as it first requires constructing a reward model to capture human preferences, and subsequently relies on reinforcement learning to tune the original unsupervised LM towards maximizing this surrogate reward while minimizing divergence from the base model. In this work, we propose a novel way to parameterize the reward model within RLHF, enabling direct recovery of the optimal policy in closed form. This insight allows us to address the classical RLHF setup with merely a simple classification loss. The approach we introduce, termed \textit{Direct Preference Optimization} (DPO), is not only more robust and efficient, but also reduces computational burden—removing the necessity for language model sampling during fine-tuning and reducing extensive hyperparameter searches. Empirical results indicate that DPO is capable of adapting LMs to align with human feedback as effectively or better than current leading techniques. In particular, DPO-based fine-tuning surpasses PPO-driven RLHF in manipulating generation sentiment, and either matches or enhances summary and single-turn dialogue response quality—all while being far easier to train and deploy.
\end{abstract}

\section{Introduction}
Large-scale unsupervised language models (LMs) trained on massive corpora have demonstrated unexpected competencies~\citep{chowdhery2022palm, brown2020language, touvron2023llama,bubeck2023sparks}. Nevertheless, these models are exposed to data authored by individuals with diverse objectives, priorities, and proficiencies. Certain intentions and expertise present in the training data are not always appropriate for models to emulate; for instance, an AI coding assistant should \textit{recognize} prevalent programming errors to provide corrections, yet when synthesizing code, it should preferentially reflect the (possibly infrequent) examples of high-caliber programming in its training distribution. Likewise, although the language model ought to be \textit{informed} about a widely held misconception—say, one believed by half the population—it should not propagate this falsehood as fact in half its responses to pertinent questions. Therefore, differentiating a model’s \emph{preferred responses and conduct} from the broader array of its \textit{knowledge and capabilities} is essential for designing AI systems that are secure, effective, and steerable \citep{ouyang2022training}. Conventional strategies primarily guide LMs towards human-preferred outputs through reinforcement learning (RL). In this work, we demonstrate that the RL-based objective commonly adopted in existing frameworks can be solved precisely using a straightforward binary cross-entropy loss, thereby streamlining the overall process of preference alignment.

Broadly, current techniques guide language models toward target behaviors by employing selected sets of human preference data indicative of responses deemed helpful and safe. This preference-alignment phase follows a preliminary unsupervised pre-training phase over extensive text data. While supervised fine-tuning on exemplary human-provided responses represents the most direct path for incorporating preferences, the prevailing and most effective methodology is reinforcement learning from human (or AI) feedback (RLHF/RLAIF; \citep{christiano2017deep,bai2022constitutional}). In RLHF, a reward model is trained using human preference data, and the LM policy is further optimized with RL to maximize rewards for desired outputs, all while restricting excessive divergence from the original model’s behavior. Though RLHF yields models capable of sophisticated dialogue and programming tasks, it introduces substantial complexity compared to standard supervised approaches, involving iterative training of several language models and integrating model sampling into the RL optimization process, which collectively result in increased computational demands.

This work presents a method for directly training language models in accordance with human preferences, bypassing the need for explicit reward modeling or reinforcement learning. We introduce \textit{{Direct Preference Optimization} (DPO)}, an algorithm that, while implicitly maximizing the same objective targeted by standard RLHF approaches (i.e., reward maximization under a KL-divergence regularization), offers a more straightforward and easily implemented solution. At a high level, DPO updates adjust the log odds in favor of preferred responses compared to less favored ones, incorporating an adaptive, example-specific importance weighting mechanism. This addition is crucial, as it avoids the model collapse that can result from directly using probability ratio-based objectives. Like prior algorithms, DPO utilizes a theoretical preference model (for example, the Bradley-Terry model; \cite{bradley1952rankanalysis}) to gauge the agreement between a reward function and observed preference data. Distinctly, rather than training a reward model based on preference loss and subsequently optimizing a policy with respect to this learned reward, DPO employs a variable transformation, expressing the preference loss directly in terms of the policy. With access to a dataset of human judgments regarding model outputs, DPO enables policy training via a straightforward binary cross entropy loss, effectively recovering an optimal policy with respect to an implicit reward function consistent with the collected preferences.

The central advance presented here is Direct Preference Optimization (DPO), an RL-free approach to preference-based language model training. Empirical evaluations demonstrate that DPO matches or exceeds the performance of current techniques—including PPO-based RLHF—in a variety of preference-driven tasks, such as sentiment adaptation, text summarization, and conversational modeling, using language models scaling up to 6B parameters.

\section{Related Work}

As self-supervised language models scale up, they demonstrate growing capability to perform tasks in a zero-shot manner \citep{radford2019language} or by leveraging only a handful of examples in few-shot prompts \citep{gpt3,megatron,chowdhery2022palm}. Nevertheless, their effectiveness on specific downstream applications and the degree to which their outputs align with user expectations can be notably enhanced through fine-tuning on collections of instructions and human-authored completions \citep{mishra-etal-2022-cross,sanh2022multitask,chung2022scaling,thoppilan2022lamda}. This process, known as `instruction-tuning', equips large language models (LLMs) with the ability to extrapolate to previously unseen instructions, thereby broadening their utility \citep{chung2022scaling}. While instruction tuning has yielded substantial advances, it is often more feasible to collect \textit{comparative} human feedback on response quality than to gather expert-curated demonstrations. Consequently, a line of research has adapted LLMs to human preferences using datasets comprised of such feedback, which has led to notable improvements in areas like translation \citep{kreutzer-etal-2018-reliability}, summarization \citep{stiennon2022learning,ziegler2020finetuning}, creative writing \citep{ziegler2020finetuning}, and adherence to instructions \citep{ouyang2022training,ramamurthy2023is}. Typically, these approaches involve first fitting a neural reward function to human preference data using frameworks such as the Bradley-Terry model \citep{bradley1952rankanalysis}, followed by refining the language model via reinforcement learning methods—commonly employing REINFORCE \citep{williams1992reinforce}, proximal policy optimization (PPO; \cite{schulman2017proximal}), or their adaptations \citep{ramamurthy2023is}. In related research, LLMs that have undergone instruction-following fine-tuning, incorporating human feedback, are used to synthesize additional preference data for specific traits like safety or harmlessness \citep{bai2022constitutional}, relying primarily on light human oversight through textual rubrics for the LLM’s evaluations. Collectively, these strategies bridge two distinct research directions: one focuses on applying reinforcement learning to train language models for diverse goals~\citep{Ranzato2015SequenceLT,paulus2018a,wu2018learning}, while the other pursues general methodologies for incorporating human preferences \citep{christiano2017deep,kupcsik2018learning}. Despite the theoretical advantages of optimizing with relative human preferences, leveraging reinforcement learning for large-scale language model fine-tuning remains a considerable practical hurdle. The present work introduces a principled alternative for optimizing based on relative preferences, obviating the need for reinforcement learning.

Research on learning policies from preferences outside the domain of language has been conducted in both the contextual bandit and reinforcement learning paradigms, leading to the development of multiple methodologies. In contextual bandit frameworks, policies are often learned based on preferences or rankings among actions instead of relying on explicit reward signals, a setting referred to as contextual dueling bandits (CDB; \cite{yue2012karmed, dudik2015contextual}). Without access to absolute reward feedback, the theoretical analysis of CDBs introduces the concept of a \textit{von Neumann winner}, defined as a policy whose expected probability of outperforming any other policy reaches at least 50\% \citep{dudik2015contextual}. Notably, in the CDB scenario, preference feedback is supplied in an online fashion, whereas learning from human preferences commonly proceeds from a static, offline dataset composed of annotated action pairs \citep{yan2022human}. In a related vein, \textit{preference-based RL} (PbRL) frameworks operate by interpreting binary preference information originating from an unknown `scoring' function, as opposed to utilizing reward signals \citep{BusaFekete2014, ruiz2023dueling}. A wide variety of algorithms for PbRL have been put forth, some of which are capable of leveraging off-policy preference data. These approaches typically first require inferring the latent scoring function—i.e., constructing a reward model—before optimizing the policy with respect to it \citep{jain2013learning, BusaFekete2014, christiano2017deep, sadigh2017active, kupcsik2018learning}. In contrast, we propose a unified policy learning method that directly updates the policy to align with expressed preferences, bypassing separate reward modeling stages.

\section{Preliminaries}\label{section:prelims}

We revisit the RLHF procedure detailed by \citeauthor{ziegler2020finetuning}, subsequently extended by works such as \citep{stiennon2022learning, bai2022training, ouyang2022training}. This methodology is generally composed of three key stages: (1) supervised fine-tuning (SFT); (2) gathering preferences and modeling the reward; and (3) reinforcement learning-based optimization.

\textbf{SFT}: The process of RLHF is typically initiated by adapting a pre-trained language model to target tasks (such as dialogue or summarization) through supervised learning on curated, high-quality datasets. This results in a model denoted as $\pi^\text{SFT}$.

\textbf{Reward Modelling Phase}: In the next stage, the SFT model is supplied with prompts $x$, and is used to generate answer pairs $(y_1, y_2) \sim \pi^\text{SFT}(y \mid x)$. Human annotators are then tasked with selecting their preferred response between the two, denoted as $y_w \succ y_l \mid x$, where $y_w$ is the favored output and $y_l$ is the less preferred one. These preferences are postulated to stem from an unobserved, underlying reward function $r^*(y, x)$ that remains inaccessible. Multiple methods exist to model such preference information, with the Bradley-Terry (BT) model \cite{bradley1952rankanalysis} being widely adopted (though generalizations such as the Plackett-Luce model \citep{plackett1975analysis, luce2012individual} are equally applicable, particularly when handling rankings over more than two answers). The BT framework describes the probability distribution over human preferences, denoted as $p^*$, as follows:

\begin{equation}\label{eq:bradley-terry}
    p^*(y_1\succ y_2 \mid x)=\frac{\exp\left(r^*(x, y_1)\right)}{\exp\left(r^*(x, y_1)\right) + \exp\left(r^*(x, y_2)\right)}.
\end{equation}

Suppose we are provided with a fixed dataset of preference pairs $\mathcal{D}=\bigl\{x^{(i)}, y_w^{(i)}, y_l^{(i)}\bigr\}_{i=1}^N$ drawn according to $p^*$. We introduce a parameterized reward function $r_{\phi}(x, y)$ and aim to fit its parameters by maximizing the likelihood of observed preferences. Formulating this as a binary classification problem leads to the following negative log-likelihood objective:

\begin{equation}\label{eq:reward_model}
    \mathcal{L}_R(r_{\phi}, \mathcal{D}) = -\mathbb{E}_{(x, y_w, y_l)\sim \mathcal{D}}\bigl[\log \sigma(r_{\phi}(x, y_w)- r_{\phi}(x, y_l))\bigr]
\end{equation}

with $\sigma$ denoting the logistic sigmoid function. For language model applications, the reward function $r_{\phi}(x, y)$ is generally initialized from the SFT model $\pi^\text{SFT}(y \mid x)$, augmented by a linear projection added atop the final transformer block, resulting in a scalar reward output \cite{ziegler2020finetuning}. To reduce reward function variance, previous work has standardized the reward values to satisfy $\mathbb{E}_{x,y\sim \mathcal{D}}\left[r_\phi(x, y)\right] = 0$ over $x$.

\textbf{RL Fine-Tuning Phase}: Within the RL fine-tuning process, this learned reward acts as a feedback mechanism for guiding language model outputs. Consistent with approaches from earlier research~\citep{jaques2017sequence, jaques2020human}, the objective is expressed as

\begin{equation}\label{eq:RL}
\max_{\pi_{\theta}}  \mathbb{E}_{x\sim \mathcal{D}, y\sim \pi_{\theta}(y \mid x)}\bigl[r_{\phi}(x, y)\bigr] - \beta\mathbb{D}_{\textrm{KL}}\bigl[\pi_{\theta}(y\mid x)\mid \mid \pi_\text{ref}(y\mid x)\bigr],
\end{equation}

where $\beta$ serves as a regularization weight controlling divergence from the reference policy $\pi_\text{ref}$, typically taken to be the initial SFT model $\pi^\text{SFT}$. In practice, the policy $\pi_\theta$ is also initialized from $\pi^\text{SFT}$. The KL penalty plays a critical role, preventing excessive shift away from data distributions where the reward model remains valid and helping maintain diverse generations by avoiding collapse onto isolated high-reward outputs. Since text generation involves discrete sampling, this training objective is not directly differentiable, necessitating reinforcement learning techniques for optimization. The widely adopted methodology \citep{ziegler2020finetuning, stiennon2022learning, bai2022training, ouyang2022training} defines the reward as ${r(x, y) = r_{\phi}(x, y) -\beta (\log \pi_{\theta}(y\mid x) - \log \pi_\text{ref}(y\mid x))}$ and applies PPO \cite{schulman2017proximal} to maximize this formulation.

\section{Direct Preference Optimization}\label{sec:DPO}

Recognizing the difficulties that reinforcement learning can present when scaling to fine-tune large language models, we seek a more streamlined policy optimization approach that works directly with preference data. Instead of the two-stage RLHF pipeline—first fitting a reward model, then applying reinforcement learning—our method is predicated on selecting a particular parameterization of the reward model that yields a closed-form optimal policy, removing the need for iterative RL updates. As described in detail below, the core idea is to exploit an analytical correspondence between reward functions and their induced optimal policies. This allows the reformulation of loss functions originally defined on reward models as loss functions over policies. In doing so, our method sidesteps learning a separate, explicit reward function, while remaining consistent with standard preference modeling frameworks, including the Bradley-Terry model. Ultimately, the policy network now implicitly serves

\textbf{Derivation of the DPO Objective.} Beginning with the standard reinforcement learning (RL) objective as outlined in Eq.~\ref{eq:RL}, applied to a generic reward function $r$, and drawing from previous research~\citep{peters2007reinforcement, peng2019advantage, korbak2022reinforcement, go2023aligning}, it can be demonstrated that the optimal policy under a KL-regularized reward maximization framework, as described in Eq.~\ref{eq:RL}, has the following closed-form:

\begin{equation}\label{eq:op_policy}
    \pi_r(y\mid x) = \frac{1}{Z(x)}\pi_\text{ref}(y\mid x)\exp\left(\frac{1}{\beta}r(x, y)\right),
\end{equation}

where the normalization constant $Z(x)$ is defined as $Z(x) =\sum_{y}\pi_\text{ref}(y\mid x)\exp\left(\frac{1}{\beta}r(x, y)\right)$. For a detailed derivation, see Appendix~\ref{app:derivation1}. Even when substituting the MLE-based estimator $r_{\phi}$ for the unknown true reward function $r^*$, estimating $Z(x)$ remains computationally prohibitive~\citep{korbak2022reinforcement, go2023aligning}, limiting the practicality of this formulation. Nonetheless, Eq.~\ref{eq:op_policy} can be manipulated to write the reward as a function of the optimal policy $\pi_r$, the reference policy $\pi_\text{ref}$, and the partition function $Z(\cdot)$. By taking the logarithm of Eq.~\ref{eq:op_policy} and rearranging, one arrives at:

\begin{equation}\label{eq:main_eq}
    r(x,y) =\beta \log \frac{\pi_r(y\mid x)}{\pi_\text{ref}(y\mid x)} + \beta \log Z(x).
\end{equation}

This reparameterization extends to both the true reward $r^*$ and its optimal policy $\pi^*$. Notably, within the Bradley-Terry model framework, what matters is only the difference in rewards across two samples; specifically, ${p^*(y_1 \succ y_2 \mid x) = \sigma(r^*(x, y_1) - r^*(x, y_2))}$. By substituting the above reparameterization from Eq.~\ref{eq:main_eq} for $r^*(x, y)$ into the preference model of Eq.~\ref{eq:bradley-terry}, the partition function terms cancel, allowing us to express the human preference probabilities solely in terms of the optimal policy $\pi^*$ and the reference $\pi_\text{ref}$. Consequently, the optimal RLHF policy $\pi^*$ is consistent with the Bradley-Terry preference structure, which leads to:

\begin{equation}\label{eq:objective}
    p^*(y_1\succ y_2 \mid x)=\frac{1}{1 + \exp\left(\beta \log \frac{\pi^*(y_2\mid x)}{\pi_\text{ref}(y_2\mid x)} - \beta \log \frac{\pi^*(y_1\mid x)}{\pi_\text{ref}(y_1\mid x)}\right)}
\end{equation}

A detailed walkthrough is provided in Appendix~\ref{app:derivation2}. Although Eq.~\ref{eq:objective} leverages the Bradley-Terry framework, parallel derivations for more general preference models, such as the Plackett-Luce models~\citep{plackett1975analysis, luce2012individual}, are available in Appendix~\ref{app:plackett_luce_models}.

Having established an expression for the likelihood of human preference data in terms of the optimal policy, rather than the reward, we can articulate a maximum likelihood estimation objective for a parametrized policy $\pi_\theta$. Similar to reward modeling as per Eq.~\ref{eq:reward_model}, this leads to the following policy objective:

\begin{equation}\label{eq:optimum_model}
    \mathcal{L}_\text{DPO}(\pi_{\theta}; \pi_\text{ref}) = -\mathbb{E}_{(x, y_w, y_l)\sim \mathcal{D}}\left[\log \sigma \left(\beta \log \frac{\pi_{\theta}(y_w\mid x)}{\pi_\text{ref

In this approach, we employ an alternative parameterization to implicitly model the reward function, so that the optimal policy corresponds directly to $\pi_\theta$. Furthermore, as this method is equivalent to fitting a reparameterized Bradley-Terry model, it benefits from certain theoretical guarantees, including consistency given appropriate assumptions about the underlying preference data distribution \cite{bong2022generalized}. We provide additional discussion on the theoretical foundations of DPO and its relationship to related methodologies in Section~\ref{sec:theory}.

\textbf{What does the DPO update accomplish?} To gain insight into the operational mechanism of DPO, it is instructive to examine the gradient of the loss $\mathcal{L}_\text{DPO}$. The derivative with respect to $\theta$ can be expressed as:

\begin{multline*}\label{eq:gradient}
    \nabla_\theta \mathcal{L}_\text{DPO}(\pi_\theta;\pi_\text{ref}) = \\ -\beta\mathbb{E}_{(x, y_w, y_l) \sim \mathcal{D}} \bigg[\underbrace{\sigma(\hat{r}_\theta(x, y_l) - \hat{r}_\theta (x, y_w))}_\text{greater emphasis when the reward estimate is inaccurate}\bigg[\underbrace{\nabla_\theta\log \pi(y_w \mid x)}_\text{promote probability of $y_w$} - \underbrace{\nabla_\theta\log\pi(y_l \mid x)}_\text{reduce probability of $y_l$}\bigg]\bigg],
\end{multline*}

Here, $\hat{r}_\theta(x, y) = \beta \log \frac{\pi_\theta(y \mid x)}{\pi_\text{ref}(y \mid x)}$ defines an implicit reward in terms of the language model $\pi_\theta$ and the reference model $\pi_\text{ref}$ (elaborated in Section~\ref{sec:theory}). Conceptually, the gradient of $\mathcal{L}_\text{DPO}$ acts to increase the probability of preferred outputs $y_w$ while reducing the likelihood assigned to less preferred ones $y_l$. Notably, the adjustment applied to each example is proportional to the extent that the implicit reward $\hat{r}_\theta$ erroneously rates the less desirable completions above the preferred ones, modulated by the scaling factor $\beta$—essentially reflecting both the misranking severity and the influence of the KL constraint. Our empirical findings underscore the significance of this adaptive weighting: omitting it (i.e., a straightforward variant without the coefficient) can lead to substantial degeneration of the language model’s outputs (see Appendix Table~\ref{tab:unlikelihood_generations}).

\textbf{DPO overview.} The standard DPO procedure involves the following steps: 1) For each prompt $x$, generate responses $y_1, y_2 \sim \pi_\text{ref}(\cdot \mid x)$ and assign human preference labels to form the offline preference dataset $\mathcal{D} = \{x^{(i)}, y_w^{(i)}, y_l^{(i)}\}_{i=1}^N$; 2) Adjust the language model $\pi_\theta$ by minimizing $\mathcal{L}_\text{DPO}$ using the given $\pi_\text{ref}$, the dataset $\mathcal{D}$, and the target parameter $\beta$. In applied scenarios, leveraging publicly available preference datasets is preferable over synthesizing new samples and collecting corresponding human judgments. Since these datasets are typically constructed using $\pi^\text{SFT}$, we default to setting $\pi_\text{ref} = \pi^\text{SFT}$ when it can be accessed. In the absence of $\pi^\text{SFT}$, we instead fit $\pi_\text{ref}$ by maximizing the likelihood of chosen completions ${(x, y_w)}$, specifically by solving ${\pi_\text{ref} = \argmax_{\pi}\mathbb{E}_{x, y_w \sim \mathcal{D}}\left[\log \pi(y_w \mid x)\right]}$. This approach is designed to counteract distributional mismatches between the intractable true reference distribution and the $\pi_\text{ref}$ utilized within DPO. Additional implementation specifics and details on hyperparameter choices are provided in Appendix~\ref{app:implementation}.

\section{Theoretical Analysis of DPO} This section elaborates on the underlying principles of the DPO method, establishes its theoretical properties, and clarifies the benefits of DPO, particularly in contrast to challenges associated with actor-critic methods in RLHF, including approaches such as PPO~\cite{schulman2017proximal}.

\label{sec:theory}

\subsection{Your Language Model Is Secretly a Reward Model} DPO achieves direct policy optimization by reframing both reward modeling and policy training as a single maximum likelihood problem, avoiding explicit reward function fitting and subsequent reinforcement learning steps. Importantly, the objective function in Eq. \ref{eq:main_eq} corresponds to a Bradley-Terry model where the reward is parameterized as $r^*(x, y) = \beta \log\frac{\pi^*_\theta(y \mid x)}{\pi_\text{ref}(y \mid x)}$, and optimizing the model $\pi_{\theta}$ aligns with the transformation of reward model optimization shown in Eq. \ref{eq:reward_model}. Here, we provide a theoretical framework for this reparameterization, demonstrating that it does not restrict the set of attainable reward models and guarantees the ability to recover the optimal policy precisely. We start by introducing an equivalence relation among reward functions.

\begin{definition}
Two reward functions $r(x, y)$ and $r'(x, y)$ are said to be equivalent if and only if there exists a function $f$ such that ${r(x, y)-r'(x, y) = f(x)}$.
\end{definition}

It is straightforward to verify that this defines an equivalence relation, which divides the set of possible reward functions into equivalence classes. This leads to the following lemmas:

\begin{lemma}\label{lemma:same_prefrence}
Within the Plackett-Luce framework, and notably for the Bradley-Terry model, any two reward functions belonging to the same equivalence class induce identical preference distributions.
\end{lemma}

\begin{lemma}\label{lemma:same_policy}
    Reward functions within a given equivalence class produce the same optimal policy in the context of the constrained RL problem.
\end{lemma}

Proofs of these statements are straightforward and can be found in Appendix \ref{app:lemma1}. The first lemma reflects a well-known issue of non-identifiability within the Plackett-Luce family of models \cite{plackett1975analysis}. Because of this ambiguity, it is generally necessary to enforce additional identifiability conditions to guarantee consistency of MLE estimates as per Eq. \ref{eq:reward_model} \cite{bong2022generalized}. The second lemma establishes that any representative reward function from a given class yields the same optimal policy; consequently, our primary objective reduces to recovering a reward function from the correct equivalence class. The following theorem is proven in Appendix~\ref{app:thm1}:

\begin{theorem}\label{thm:main}
    Provided mild assumptions, every equivalence class of reward functions compatible with Plackett-Luce (including Bradley-Terry) models can be written as ${r(x, y) = \beta \log \frac{\pi(y\mid x)}{\pi_\text{ref}(y\mid x)}}$ for some choice of model $\pi(y\mid x)$ and a fixed reference $\pi_\text{ref}(y \mid x)$.
\end{theorem}

\begin{sproof}
    Take any reward function $r(x, y)$ which gives rise to an associated optimal model $\pi_r(y \mid x)$ via Eq. \ref{eq:op_policy}. We wish to show that there exists a representative of the equivalence class of $r$ that admits the reparameterization stated above. Consider the mapping $f$ defined as 
\begin{equation}
    f(r; \pi_\text{ref}, \beta)(x, y) = r(x, y) - \beta\log\sum_{y}\pi_\text{ref}(y\mid x)\exp\left(\frac{1}{\beta}r(x, y)\right)
\end{equation}
The transformation $f$ effectively centers the reward function by subtracting the log-partition function with respect to $\pi_\text{ref}$, which depends solely on $x$. Thus, $f(r; \pi_\text{ref}, \beta)(x, y)$ remains within the equivalence class determined by $r(x, y)$. Substituting $r$ using the right side of Eq.~\ref{eq:main_eq}, which is valid for any $r$, yields $f(r; \pi_\text{ref}, \beta)(x, y) = \beta \log \frac{\pi_r(y\mid x)}{\pi_\text{ref}(y\mid x)}$. Therefore, $f$ produces an explicit representative of the equivalence class of $r$ in the requisite form, and thus the proposed reparameterization entails no loss of generality in modeling the reward.
\end{sproof}

An alternative interpretation of Theorem~\ref{thm:main} is that it pinpoints the specific reward function, from each equivalence class, that is selected via the DPO reparameterization. More precisely, this reward function fulfills the following condition:

\begin{equation}\label{eq:lag_p}
     \sum_{y}\underbrace{\pi_\text{ref}(y\mid x)\exp\left(\frac{1}{\beta}r(x, y)\right)}_{=\pi(y\mid x)\text{, using Thm.~\ref{thm:main} reparam.}} = 1,
\end{equation}

This condition guarantees that $\pi(y\mid x)$ constitutes a valid probability distribution—specifically, its values are non-negative and the probabilities sum to unity. Additionally, referencing Eq.~\ref{eq:op_policy}, we recognize that Eq.~\ref{eq:lag_p} corresponds to the partition function associated with the optimal policy derived from the reward function $r(x, y)$. The central innovation of the DPO method lies in its ability to enforce constraints on the generally underdetermined Plackett-Luce class (including, in particular, the Bradley-Terry models) of preference models. By imposing these constraints, the algorithm ensures the set of expressible reward models remains intact while also rendering the optimal policy in Eq.~\ref{eq:op_policy} analytically tractable for any given prompt $x$.

\subsection{Instability of Actor-Critic Algorithms} 

This framework is equally valuable in identifying sources of instability in traditional actor-critic methods—such as PPO—applied to RLHF. Adhering to the typical RLHF sequence and emphasizing the RL fine-tuning phase described in Section \ref{section:prelims}, we establish links to the control as inference paradigm \cite{levine2018reinforcement} formulated for the constrained RL objective in \ref{eq:RL}. Suppose we utilize a parameterized distribution $\pi_{\theta}(y\mid x)$, aiming to minimize $\mathbb{D}_{\text{KL}}[\pi_{\theta}(y|x) \mid \mid \pi^*(y\mid x)]$, where $\pi^*$—the optimal policy—originates from Eq.~\ref{eq:optimum_model} and is driven by the reward function $r_{\phi}(y, x)$. Upon rearrangement, this objective reduces to:

\begin{equation}\label{eq:AC}
    \max_{\pi_{\theta}}\mathbb{E}_{\pi_{\theta}(y\mid x)}\bigg[\underbrace{r_{\phi}(x, y) -\beta\log\sum_{y}\pi_\text{ref}(y\mid x)\exp\left(\frac{1}{\beta}r_{\phi}(x, y)\right)}_{f(r_{\phi}, \pi_\text{ref}, \beta)} - \underbrace{\beta\log\frac{\pi_{\theta}(y\mid x)}{\pi_\text{ref}(y\mid x)}}_{\text{KL}}\bigg]
\end{equation}

This objective has also been utilized in earlier studies 
\citep{ziegler2020finetuning, stiennon2022learning, bai2022training, ouyang2022training}, which employed a DPO-equivalent reward formulation for the reward class $r_{\phi}$. Within this context, the normalization component in $f(r_{\phi}, \pi_\text{ref}, \beta)$ can be viewed as representing the soft value function associated with the reference policy $\pi_\text{ref}$. Although this normalization term does not alter the optimum of the objective, omitting it may result in a policy gradient with elevated variance, thereby destabilizing learning. One way to handle this normalization is by introducing a learned value function, yet optimizing such a function is itself a challenging problem. An alternative strategy, adopted in prior research, is to normalize the rewards with respect to a human completion baseline, which essentially estimates the normalization factor via a single-sample Monte-Carlo approach. In contrast, the DPO reparameterization leads to a reward function formulation that eliminates the need for any baseline corrections. 

\section{Experiments}
Here, we provide an empirical assessment of DPO's capacity to learn policies directly from human preference data. To begin, we consider a controlled text-generation environment, focusing on the question: relative to standard preference optimization algorithms such as PPO, how well does DPO balance reward maximization against minimization of KL-divergence with the reference policy? Subsequently, we expand our evaluation to include larger models and more complex RLHF scenarios, covering applications such as summarization and dialogue. Remarkably, we observe that DPO achieves comparable or superior performance to robust methods like RLHF with PPO, and outperforms returning the highest-reward trajectory from $N$ sampled outputs—often with minimal or no hyperparameter adjustment. Prior to reporting these empirical findings, we outline the experimental protocols; further specifics can be found in Appendix~\ref{app:exp_details}.

\textbf{Tasks.} We conduct our study on three distinct open-domain text generation challenges. In all cases, learning algorithms infer a policy using a set of preferences $\mathcal{D}=\bigl\{x^{(i)}, y_w^{(i)}, y_l^{(i)}\bigr\}_{i=1}^N$. The first task, \textbf{controlled sentiment generation}, presents $x$ as an initial segment of a movie review from the IMDb dataset \cite{maas-EtAl:2011:ACL-HLT2011}; the objective for the policy is to produce a continuation $y$ expressing positive sentiment. To facilitate systematic assessment, we automatically construct pairs of generations for comparison via a pre-trained sentiment classifier, ensuring that $p(\text{positive}\mid x,y_w)>p(\text{positive}\mid x,y_l)$. For SFT, we adapt GPT-2-large via fine-tuning on training-set reviews from IMDb until training stabilizes (additional methodological details are presented in App~\ref{app:sentiment_details}). 

In the \textbf{summarization} task, $x$ refers to a Reddit forum post, with the policy required to generate a concise summary $y$ that captures the essential points. As in earlier research, our experiments rely on the Reddit TL;DR summarization dataset \citep{volske-etal-2017-tl} combined with human-labeled preference data collected by \citeauthor{stiennon2022learning}. For SFT, we employ a model trained on human-authored Reddit summaries\footnote{\url{https://huggingface.co/CarperAI/openai_summarize_tldr_sft}} using the TRLX \citep{leandro_von_werra_2023_7790115} system for RLHF. Notably, the human preference annotations were generated from samples of a different, but comparably-trained, SFT model. 

The third task, \textbf{single-turn dialogue}, frames $x$ as a user prompt—ranging from technical inquiries to personal questions. Here, the policy is tasked with generating a response $y$ that is both helpful and engaging. For this, we utilize the Anthropic Helpful and Harmless dialogue dataset \citep{bai2022training}, comprising 170,000 human-assistant conversations. Each conversation concludes with two model-generated replies and a human preference indicator for the superior response. As no pre-existing SFT model is available for this dataset, we develop our SFT baseline by fine-tuning an out-of-the-box language model using only the responses favored by human raters.

\textbf{Evaluation.} Our experimental framework adopts two primary evaluation strategies. To directly assess each algorithm's ability to optimize for constrained reward maximization, we measure its performance in controlled sentiment generation by examining the set of achievable reward and KL-divergence trade-offs relative to a reference policy. This evaluation is possible in this context because the reward function—a sentiment classifier—serves as an accessible oracle. Conversely, in practical scenarios where such a ground-truth reward function is unavailable, we resort to measuring each algorithm's \textit{win rate} compared to a baseline policy. For these real-world summarization and single-turn dialogue tasks, GPT-4 serves as a stand-in for human evaluators, judging both summary quality and the helpfulness of dialogue responses. Specifically, the baseline in summarization corresponds to the reference summaries from the test data, while for dialogue it is the human-preferred test response. While some research indicates language models can outperform existing automatic metrics as evaluators \citep{Chen2023ExploringTU}, we validate our choice of GPT-4 for assessment through a human study presented in Sec.~\ref{sec:human-judgments}. Our findings demonstrate that GPT-4’s ratings are strongly aligned with those of human annotators, and in many cases, humans concur with GPT-4 at rates equal to or higher than their agreement with each other.

\textbf{Methods.} In our evaluation, alongside DPO, we consider a variety of established techniques for aligning language model outputs with human preference. As baselines, we employ zero-shot prompting with \textbf{GPT-J} \citep{gpt-j} for the summarization experiments and 2-shot prompting with \textbf{Pythia-2.8B} \citep{biderman2023pythia} for dialogue generation. We further include models trained by \textbf{SFT} as well as \textbf{Preferred-FT}, which is generated by fine-tuning via supervised learning using preferred completions $y_w$; these preferred examples are sourced from either the SFT model (for sentiment and summarization tasks) or a general-purpose language model (for dialogue). Another supervised-like baseline, \textbf{Unlikelihood}~\citep{welleck2019neural}, modifies the training objective by encouraging higher probability for $y_w$ while penalizing likelihood for $y_l$, controlled by an adjustable coefficient $\alpha\in[0,1]$. Additionally, we examine \textbf{PPO} \citep{schulman2017proximal}, in which rewards are estimated from preference annotations, and its oracle variant, \textbf{PPO-GT}, that leverages the true reward signal in the controlled sentiment setup. For sentiment experiments, we assess both a standard \cite{leandro_von_werra_2023_7790115} implementation of PPO-GT and a modified version that normalizes the rewards and tunes hyperparameters for greater effectiveness (similar modifications apply to PPO with learned rewards). Lastly, we include the \textbf{Best of $N$} strategy, which involves sampling $N$ candidate outputs from the SFT (or Preferred-FT for dialogue), and returning the response that achieves the highest score as per the preference-learned reward model. Despite its strong empirical results, this method separates reward modeling from PPO optimization but entails substantial computational costs, as it demands generating $N$ completions per prompt during evaluation.

\subsection{How well can DPO optimize the RLHF objective?}

In conventional RLHF approaches, the KL-constrained reward maximization objective serves to balance obtaining higher rewards with ensuring that the learned policy does not stray significantly from the reference policy. Thus, any comparison between methods should jointly consider both the attained reward and the KL divergence; optimizing for slightly higher rewards at the cost of substantially increasing KL is not inherently advantageous. Figure~\ref{fig:frontier-tldr-main} presents the tradeoff between reward and KL for a range of methods within the sentiment scenario. For each method, we conduct several experiments, varying the policy conservativeness parameter each time (target KL values of $\{3,6,9,12\}$ for PPO, $\beta$ values in $\{0.05,0.1,1,5\}$, $\alpha$ values from $\{0.05,0.1,0.5,1\}$ for unlikelihood, and different random seeds for preferred-FT). In total, this grid search comprises 22 experiments. At intervals of every 100 training steps, continuing until convergence, we assess each trained policy using a designated set of test prompts, recording both the mean reward according to the true reward function and the mean sequence-level KL\footnote{Here, the sequence-level KL refers to the aggregate of KL divergences across timesteps.} with respect to the reference policy, $\text{KL}\left(\pi\mid \mid \pi_\text{ref}\right)$. The results indicate that DPO yields a Pareto frontier that is markedly superior, attaining greater rewards while maintaining relatively low KL divergence. This observation is especially striking for several reasons. First, while both DPO and PPO are designed to optimize an equivalent objective, DPO demonstrates greater efficiency; the DPO reward-KL Pareto curve consistently surpasses that of PPO. Second, DPO continues to deliver a superior tradeoff compared to PPO, \emph{even in scenarios where PPO has access to the ground truth reward} (PPO-GT).

\subsection{Can DPO scale to real preference datasets?}
\label{sec:dpo-real-datasets}
We proceed to investigate the fine-tuning capabilities of DPO on tasks involving summarization and single-turn dialogue. In summarization, it is known that standard automated metrics like ROUGE often show limited alignment with actual human preferences~\citep{stiennon2022learning}. Previous research has indicated that tuning language models with PPO using human feedback can result in more effective summaries. To assess the performance of various approaches, we sample completions from the test portion of the TL;DR summarization dataset and calculate the average win rate when compared to reference completions in this test set. For each method, completions are drawn at sampling temperatures ranging from 0.0 up to 1.0, with corresponding win rates visualized in Figure~\ref{fig:frontier-tldr-main} (right). DPO, PPO, and Preferred-FT are all fine-tuned from an identical GPT-J SFT model\footnote{\url{https://huggingface.co/CarperAI/openai_summarize_tldr_sft}}. Our results indicate that DPO achieves a win rate near 61\% at a sampling temperature of 0.0, surpassing PPO, which reaches approximately 57\% at its best temperature setting of 0.0. Additionally, DPO obtains a superior peak win rate relative to the optimal $N$ baseline. It is important to highlight that DPO's $\beta$ hyperparameter was not extensively tuned in these experiments, implying that actual performance may be understated. DPO also demonstrates greater resilience to variations in sampling temperature in contrast to PPO, whose performance can decline to that of the unmodified GPT-J model as temperature increases. The Preferred-FT method yields negligible gains over the SFT model baseline. In Section~\ref{sec:human-judgments}, we present human evaluation results directly comparing DPO and PPO, finding that DPO outputs sampled at temperature 0.25 are chosen 58\% of the time over PPO outputs at the same temperature.

For the single-turn dialogue task, we conduct experiments on a portion of the test split from the Anthropic HH dataset \citep{bai2022training}, specifically focusing on examples that involve only one human-assistant exchange. To assess the different approaches, GPT-4 evaluations reference the preferred completions in the test set to calculate the win rates for each method. Due to the absence of a standardized SFT model for this scenario, our workflow initiates with a pre-trained Pythia-2.8B model. We apply Preferred-FT to finetune a reference model on the preferred completions, ensuring that generated completions align with the model’s distribution, and subsequently proceed to DPO training. Our analysis also incorporates a comparison against the top-performing out of 128 Preferred-FT completions (noting that the Best of $N$ benchmark reaches a performance plateau at 128 samples for this dataset, as demonstrated in Appendix Figure~\ref{fig:best-of-n}), as well as a 2-shot prompted version of the base Pythia-2.8B model. We observe that DPO matches or surpasses these baselines when employing the most favorable temperature settings for each approach. Additionally, we evaluate an RLHF model trained on the Anthropic HH dataset using PPO\footnote{\url{https://huggingface.co/reciprocate/ppo_hh_pythia-6B}} provided by a reputable source\footnote{\url{https://github.com/CarperAI/trlx/tree/main/examples/hh}}, but despite extensive prompt and temperature tuning, this model fails to outperform the original Pythia-2.8B. Drawing on our TL;DR findings and acknowledging that both techniques maximize the same reward signal, we regard the Best of 128 results as a practical approximation for PPO-based performance. Taken together, our findings indicate that DPO uniquely offers a computationally efficient method to exceed the preferred completions within the Anthropic HH dataset, and achieves outcomes on par with or superior to the more resource-intensive Best of 128 approach. As depicted in Figure~\ref{fig:dialogue-main}, DPO attains peak performance in a relatively small number of training steps.

\subsection{Generalization to a new input distribution}

To better assess how PPO and DPO perform when faced with distributional shifts, we analyze the generalization abilities of their policies—trained on the Reddit TL;DR summarization task—by evaluating them on a separate dataset consisting of news articles from the CNN/DailyMail test split \citep{nallapati-etal-2016-abstractive}. The evaluation employs the optimal sampling temperatures identified in the original TL;DR experiments (0 and 0.25), and the outcomes are detailed in Table~\ref{tab:ood}. We used GPT-4’s win rate metric by comparing model outputs to reference summaries, applying the same GPT-4 (C) prompt as in the Reddit TL;DR setup but substituting “forum post” with “news article”. On this unseen data distribution, the DPO policy maintains a strong lead over PPO. These findings preliminarily suggest that DPO is capable of generalizing to novel distributions as well as PPO, despite not utilizing the additional unlabeled Reddit TL;DR data that PPO does.

\subsection{Validating GPT-4 judgments with human judgments}
\label{sec:human-judgments}
A human subject study was undertaken to examine the accuracy of GPT-4 as an evaluator, leveraging results from the TL;DR summarization trials with two variants of GPT-4 prompts. The \textbf{GPT-4 (S)} (“simple”) prompt requests GPT-4 to judge which summary better conveys the essential information in a given post. The \textbf{GPT-4 (C)} (“concise”) prompt includes an additional criterion, asking which summary is also more concise. The latter prompt was investigated because, under the “simple” prompt, GPT-4 tends to favor summaries that are lengthier and more redundant compared to human preferences. The full text of the prompts can be found in Appendix~\ref{app:prompts}. We perform three pairwise evaluations—covering the strongest-performing method (DPO, temp. 0.25), the weakest (PPO, temp. 1.0), and an intermediate (SFT, temp. 0.25)—each against PPO with greedy sampling at its best-performing temperature, ensuring coverage of a broad spectrum of generation qualities. Across both prompts, GPT-4’s preferences largely align with human raters to an extent comparable with the agreement rates observed among humans themselves (for DPO and PPO-1, multiple human ratings were collected; this was limited by rater availability). In general, win rates from the \textbf{GPT-4 (C)} prompt are more closely matched to human assessments, so this prompt was adopted for main result reporting in Section~\ref{sec:dpo-real-datasets}. Further information regarding the human evaluation protocol—including the rater interface and a roster of participating human evaluators—can be found in Appendix~\ref{app:human-study}.

\section{Discussion}
Preference-based learning stands out as an effective and scalable approach for developing language models that are both capable and aligned with human intentions. In this work, we present DPO, a straightforward framework for leveraging preferences in language model training, sidestepping the need for reinforcement learning techniques. Rather than adapting preference learning into a conventional RL context to fit standard RL solutions, DPO establishes a direct correspondence between language model policies and reward functions. This direct approach enables preference-driven model training using a simple cross-entropy objective, entirely avoiding the complexities of reinforcement learning while maintaining expressive power. Remarkably, DPO achieves comparable or superior results to popular RLHF methods, such as those built on PPO, with minimal hyperparameter adjustment. As a result, DPO lowers the practical barriers associated with preference-driven training of language models.

\textbf{Limitations \& Future Work.} Our findings highlight a number of compelling directions for further investigation. For instance, how well do DPO-trained policies generalize to out-of-distribution data compared to those trained via explicit reward functions? Preliminary evidence suggests that DPO's generalization capabilities are on par with PPO-based approaches, yet more detailed exploration is required. Another open question concerns the efficacy of leveraging self-labeled data via the DPO policy to utilize unlabeled prompts effectively. Furthermore, the dynamics of reward over-optimization in the context of direct preference optimization warrant closer scrutiny—such as whether the modest decline in performance observed in Figure~\ref{fig:dialogue-main}-right exemplifies this phenomenon. Though our experiments have reached models with up to 6B parameters, future work may investigate DPO's scalability to state-of-the-art models with much larger parameter counts. With respect to evaluation, we observed that GPT-4-derived win rates are sensitive to prompt variation; determining optimal procedures for extracting reliable judgments from automated systems remains an open research problem. Ultimately, DPO's utility could extend far beyond language modeling from human preferences, offering promising opportunities for training generative models across a variety of domains.